\paragraph{6.2.3. A Generic Connection from Secret Sharing Schemes.}
We now formally state a generic reduction that reduces the problem of estimating the weight of access structures to $\mathsf{GapMINcKT}$.
\paragraph{Theorem 6.14.}
Let $\mathcal A=\{\mathcal A_\phi\}_{\phi\in\{0,1\}^*}$ be a family of access structures that admits efficient secret sharing schemes. Let $\tau:\mathbb N^3\to\mathbb N$ be a function such that $t\leq\tau(n,m,t)\leq n^c\cdot m^{1-1/c}\cdot t^c$ for some constant $c>1$. Then, there exists a randomized polynomial-time algorithm $R$ that takes $\phi\in\{0,1\}^*$ as input and outputs $x,y\in\{0,1\}^*$ and $t,\rho\in\mathbb N$ such that
\[\frac\rho{8c}\cdot w(\mathcal A_\phi)+\log t'\leq K^{t'}(x\mid y)\leq K^t(x\mid y)\leq\rho\cdot w(\mathcal A_\phi)-\operatorname{cd}^t(y)\]
holds with probability at least $1-o(1)$ over the internal randomness of $R$, where $t':=\tau(|x|,|y|,t)$.
It is easy to observe that this reduction implies $\mathsf{NP}$-hardness of $\mathsf{GapMINcKT}$.
\begin{proof}[Proof of Theorem 6.6]
Let $\mathcal A=\{\mathcal A_\phi\}_\phi$ be the family of access structures for monotone formulas given in Lemma 6.11. Let $R$ be the reduction of Theorem 6.14. Consider a reduction $R'$ that reduces an instance $(\phi,s)$ of $\mathsf{Gap}_g\mathsf{MMSA}$ to an instance $(x,y,1^t,1^{s'})$ of $\mathsf{Gap}_{\tau}\mathsf{MINcKT}$ such that $(x,y,t,\rho):=R(\phi)$ and $s':=s\cdot\rho$.
Let $(\Pi_{\mathrm{Yes}},\Pi_{\mathrm{No}}):=\mathsf{Gap}_g\mathsf{MMSA}$. We claim below that the reduction $R'$ reduces $\mathsf{Gap}_g\mathsf{MMSA}$ to $\mathsf{Gap}_{\tau,\sigma}\mathsf{MINcKT}$ for some $\sigma$, which implies $\mathsf{NP}$-hardness of $\mathsf{Gap}_{\tau,\sigma}\mathsf{MINcKT}$ by Lemma 6.13. If $(\phi,s)\in\Pi_{\mathrm{Yes}}$, then $w(\mathcal A_\phi)=\mathsf{MMSA}(\phi)\leq s$. By the properties of $R$, with high probability, it holds that
\[K^t(x\mid y)\leq\rho\cdot w(\mathcal A_\phi)-\operatorname{cd}^t(y)\leq s'-\operatorname{cd}^{t,t'}(y),\]
which implies that $(x,y,1^t,1^{s'})$ is a Yes instance of $\mathsf{Gap}_{\tau,\sigma}\mathsf{MINcKT}$. If $(\phi,s)\in\Pi_{\mathrm{No}}$, then $w(\mathcal A_\phi)=\mathsf{MMSA}(\phi)>g(|\phi|)\cdot s$. Also by the properties of $R$, with high probability, it holds that
\[K^{t'}(x\mid y)\geq\frac\rho{8c}\cdot w(\mathcal A_\phi)+\log t'>\frac{g(|\phi|)}{8c}\cdot s'+\log t',\]
which implies that $(x,y,1^t,1^{s'})$ is a No instance of $\mathsf{Gap}_{\tau,\sigma}\mathsf{MINcKT}$ for some $\sigma(|x|):=g(|\phi|)/(8c)\geq|\phi|^{1/(\log\log|\phi|)^{O(1)}}\geq|x|^{1/(\log\log|x|)^{O(1)}}$, where the last inequality follows from the fact that $|\phi|\geq|x|^{\Omega(1)}$.
\end{proof}
It remains to prove Theorem 6.14. To show that the additive error $\operatorname{cd}^t(y)$ is relatively small, we use the fact that for $D$ independent random samples $y^1,\ldots,y^D$ from a distribution $D$ samplable by polynomial-size circuits, the amortized time-bounded Kolmogorov complexity of $y^1,\ldots,y^D$ approaches the entropy of $D$ asymptotically.
We are now ready to prove Theorem 6.14.
\begin{proof}[Proof of Theorem 6.14]
Fix an input $\phi\in\{0,1\}^*$. Let $n$ be the number of parties in the access structure $\mathcal A_\phi$. Let $\lambda=O(\log|\phi|)$, $t$, and $\ell$ be parameters chosen later.
We first define a family $D=\{D_\phi\}_{\phi\in\{0,1\}^*}$ of distributions by using the following sampling procedure: Choose $x\sim\{0,1\}^\ell$ uniformly at random. Let $(s_1,\ldots,s_n):=\mathsf{Share}(\phi,x)$ and let $m$ be the length of each share; i.e., $m:=|s_i|$. Pick $k_1,\ldots,k_n\sim\{0,1\}^\lambda$ randomly. We define a string $y\in\{0,1\}^{2^\lambda\cdot m}$, which we identify with a function $y:\{0,1\}^\lambda\to\{0,1\}^m$: For every $q\in\{0,1\}^\lambda$, let $y(q):=s_i$ if $q=k_i$ for some $i\in[n]$ and let $y(q):=0^m$ otherwise. Note that $y$ is not well defined if $(k_i)_{i\in[n]}$ is not pairwise distinct; however, we will show that $y$ is well defined with high probability. The output of the sampling procedure is defined as $(x,y,k,s)$, where $k:=(k_1,\ldots,k_n)$ and $s:=(s_1,\ldots,s_n)$.
The algorithm $R$ operates as follows: For a given input $\phi\in\{0,1\}^*$, pick $D$ independent samples $(x^1,y^1,k^1,s^1),\ldots,(x^D,y^D,k^D,s^D)$ from $D_\phi$. Let $x:=(x^1,\ldots,x^D)$, $y:=(y^1,\ldots,y^D)$, and $\rho:=2\lambda D$. The output of $R$ is defined to be $(x,y,t,\rho)$.
Below, we prove the correctness of $R$ using a sequence of claims. Let $k:=(k_i^d:i\in[n],d\in[D])$ and let $s:=(s_i^d:i\in[n],d\in[D])$. We first observe that $k$ is Kolmogorov-random with high probability:
\paragraph{Claim 6.17.}
With probability at least $1-o(1)$, it holds that
\[K(k\mid s)\geq nD\lambda-\log n.\tag{16}\]
\begin{proof}
Since $k=(k_i^d)_{i\in[n],d\in[D]}$ is uniformly distributed over $(\{0,1\}^\lambda)^{nD}$ and independent of $s$, the claim follows from Fact 3.2.
\end{proof}
\paragraph{Claim 6.18.}
If Equation (16) holds, then $k=(k_i^d)_{i\in[n],d\in[D]}$ is pairwise distinct and hence $y$ is well defined.
\begin{proof}
If there exists $(i,d)\neq(i',d')$ such that $k_i^d=k_{i'}^{d'}$, then $k$ can be described by a program of size $nD\lambda-\lambda+O(\log nD)$. Choosing a large enough $\lambda=O(\log|\phi|)$, this implies that $K(k)\leq nD\lambda-\lambda+O(\log nD)<nD\lambda-\log n$, which contradicts Equation (16).
\end{proof}
Below, we assume that Equation (16) holds. We prove an upper bound of $K^t(x\mid y)$.
\paragraph{Claim 6.19.}
$K^t(x\mid y)\leq\rho\cdot w(\mathcal A_\phi)-\operatorname{cd}^t(y)$ holds with probability at least $1-o(1)$.
\begin{proof}
Let $T\in\mathcal A_\phi$ be a minimum authorized set of parties such that $|T|=w(\mathcal A_\phi)$. We present an oracle program $M^y$ that takes an index $d$ and outputs $x^d\in\{0,1\}^\ell$ as follows: $M^y$ takes $d\in[D]$ as input and $T$, $\{k_i^d:i\in T,d\in[D]\}$ and $\phi$ as hard-wired input, computes $\{s_i^d:i\in T\}=\{y^d(k_i^d):i\in T\}$ by making queries to the oracle $y$, and outputs $\mathsf{Rec}(\phi,s_T^d)$, which is equal to $x^d$ by the correctness of a secret sharing scheme and the assumption that $T\in\mathcal A_\phi$. The size of the oracle machine $M$ is at most
\[O(|T|\cdot\log n)+|T|\cdot D\cdot\lambda+|\phi|+O(\log D)\leq2\lambda D\cdot w(\mathcal A_\phi)-D^{1-\delta},\]
where $\delta>0$ is the constant from Corollary 6.16 and this inequality follows by choosing sufficiently large $D\geq\operatorname{poly}(|\phi|)$. Applying Corollary 6.16 to the distribution of $y$, we obtain $\operatorname{cd}^t(y)\leq D^{1-\delta}$ with probability at least $1-o(1)$; thus, $|M|\leq\rho\cdot w(\mathcal A_\phi)-\operatorname{cd}^t(y)$. The running time of $M$ is at most $\operatorname{poly}(|\phi|)$, which is independent of $\lambda$. We choose $t$ so that this running time is at most $t$.
\end{proof}
The remainder of the proof is devoted to proving the lower bound.
\paragraph{Claim 6.20.}
Let $\theta:=\frac\rho{8c}\cdot w(\mathcal A_\phi)+\log t'$. Then, $K^{t'}(x\mid y)\geq\theta$ with probability at least $1-o(1)$.
We first clarify the condition of random variables under which the claim holds. Let $s_{[n]}^{[D]\setminus\{d\}}:=(s_i^{d'}:i\in[n],d'\in[D]\setminus\{d\})$.
\paragraph{Claim 6.21.}
With probability at least $1-o(1)$, it holds that for every unauthorized set $T\notin\mathcal A_\phi$ of parties and for every $d\in[D]$,
\[K(x^d\mid s_T^d,s_{[n]}^{[D]\setminus\{d\}},k)\geq\ell-n-2\log D.\tag{17}\]
\begin{proof}
Fix $d\in[D]$. Since the random variable $x^d$ is independent of $s_{[n]}^{[D]\setminus\{d\}}$ and $k$, by Lemma 6.9,
\[\min_{T\notin\mathcal A_\phi}K(x^d\mid s_T^d,s_{[n]}^{[D]\setminus\{d\}},k)\geq\ell-n-2\log D\]
holds with probability at least $1-1/D^2$. The claim follows by taking a union bound over all $d\in[D]$.
\end{proof}
In what follows, we assume Equations (16) and (17) and prove $K^{t'}(x\mid y)\geq\theta$, which will complete the proof of Claim 6.20. Assume, by way of contradiction, that $K^{t'}(x\mid y)<\theta$. Let $M\in\{0,1\}^*$ be the description of an oracle program such that $|M|\leq\theta+O(\log\ell)$ and $M^y(d)$ outputs $x^d\in\{0,1\}^\ell$ in time $t'':=t'\ell$ for every $d\in[D]$; the machine $M$ can be constructed from a program that witnesses $K^{t'}(x\mid y)<\theta$. For each $d\in[D]$, let $T(M,d)$ be the set of indices $i\in[n]$ such that some bit of $y^d(k_i^d)\in\{0,1\}^m$ is queried during the computation of $M^y(d)$.
\paragraph{Claim 6.22.}
Let $M$ be an oracle program that runs in time $t''$ and let $\alpha:=\sum_{d=1}^D|T(M,d)|$. Then,
\[K(k\mid s)\leq|M|+(nD-\alpha)\cdot\lambda+\alpha\cdot(\log t''+\log nD)+O(\log nD).\]
\begin{proof}
By the definition of $T(M,d)$, for every $i\in T(M,d)$, there exists a time step $t_{i,d}\in[t'']$ such that $M^y$ makes a query $k_i^d$ to the oracle $y$ on input $d\in[D]$. To describe $k$ given $s$, consider the following program $M'$: $M'$ takes
\[\{(i,d,t_{i,d})\in[n]\times[D]\times[t'']:i\in T(M,d)\},\]
\[\{k_i^d\in\{0,1\}^\lambda:(i,d)\in[n]\times[D],\ i\notin T(M,d)\},\]
and $M$ as input and simulates $M^y$ on input $d$ for every $d\in[D]$. At the time step $t_{i,d}$ of the simulation of $M^y(d)$ for some $i\in T(M,d)$, $M'$ reads the query $k_i^d$ that $M^y(d)$ makes to the oracle $y^d$ and answers the query with $s_i^d$. $M'$ continues the simulation until the time step $t''$, at which point $M'$ knows $k_i^d$ for every $i\in T(M,d)$. Finally, $M'$ outputs $k$. The input of $M'$ can be encoded as a string of length $\alpha\cdot(\log t''+\log nD)+(nD-\alpha)\cdot\lambda+|M|+O(\log nD)$.
\end{proof}
It follows from Claim 6.22 and Equation (16) that
\[nD\lambda-\log n\leq|M|+(nD-\alpha)\cdot\lambda+\alpha\cdot(\log t''+\log nD)+O(\log nD),\]
which can be simplified to
\[(\lambda-\log t''-\log nD)\cdot\alpha\leq|M|+O(\log nD).\tag{18}\]
Since $t''/\ell=t'=\tau(|x|,|y|,t)\leq|x|^c\cdot|y|^{1-1/c}\cdot t^c\leq2^{(1-1/c)\cdot\lambda}\cdot(D\cdot|\phi|\cdot\ell\cdot m)^{O(c)}$, we may choose $\lambda=O(\log|\phi|)$ large enough so that $t''\leq2^{(1-1/2c)\cdot\lambda-\log nD}$. Then, by Equation (18), we obtain
\begin{align*}
\alpha&\leq\frac{2c}{\lambda}\cdot(|M|+O(\log nD))\\
&<\frac{4c}{\lambda}\cdot(\theta-\log t')=\frac\rho{2\lambda}\cdot w(\mathcal A_\phi)=D\cdot w(\mathcal A_\phi),
\end{align*}
where the second inequality follows from $|M|+O(\log nD)\leq\theta+O(\log nD\ell t')<2\theta$.\footnote{We may assume without loss of generality that $w(\mathcal A_\phi)\geq1$; then, we have $\theta\geq\Omega(\rho)\geq\Omega(D)$, which can be assumed to be larger than logarithmic terms.} It follows that there exists $d\in[D]$ such that $|T(M,d)|\leq\alpha/D<w(\mathcal A_\phi)$, which implies that
\[T(M,d)\notin\mathcal A_\phi.\]
By Equation (17), we obtain
\[K(x^d\mid s_{T(M,d)}^d,s_{[n]}^{[D]\setminus\{d\}},k)\geq\ell-n-2\log D.\]
However, this contradicts the following claim.
\paragraph{Claim 6.23.}
Let $M$ be an oracle program such that $M^y(d)$ outputs $x^d$ for every $d\in[D]$. Then, for every $d\in[D]$, it holds that
\[K(x^d\mid s_{T(M,d)}^d,s_{[n]}^{[D]\setminus\{d\}},k)\leq|M|+O(\log D).\]
\begin{proof}
Since $M^y(d)$ outputs $x^d$ without making any query in $\{k_i^d:i\in[n]\setminus T(M,d)\}$, we can define another oracle $z=(z^1,\ldots,z^D)$ such that $z^{d'}:=y^{d'}$ for every $d'\in[D]\setminus\{d\}$ and $z^d(k_i^d):=s_i^d$ for every $i\in T(M,d)$ and $z^d(k):=0^m$ otherwise, so that there is no difference between $y$ and $z$ on inputs queried by $M^y(d)$; therefore, we obtain $x^d=M^y(d)=M^z(d)$. The claim follows by observing that the oracle $z$ can be constructed from $M$, $d\in[D]$, $s_{T(M,d)}^d$, $s_{[n]}^{[D]\setminus\{d\}}$, and $k$.
\end{proof}
We conclude that $\ell-n-2\log D\leq|M|+O(\log D)\leq O(\theta)\leq O(n\lambda D)$, which is a contradiction by letting $\ell\gg O(n\lambda D)$. This completes the proof of Claim 6.20.
\end{proof}
\paragraph{Remark 6.24.}
It is not hard to observe that our NP-hardness hardness reduction also proves the NP-hardness of $\mathsf{McKTP}$ [4] because the upper bound of Claim 6.19 also holds for $KT(x\mid y)$.
